\documentclass{article}
\usepackage{amsmath,amssymb,amsthm}
\newtheorem{theorem}{Theorem}
\newtheorem{lemma}{Lemma}
\begin{document}
\section*{A043: genuine disjoint primes}
All gates are Boolean sign functions. Remove irrelevant coordinates first.
An essential gate $f$ on $n\ge2$ inputs is prime if no block
$B\subset[n]$ with $2\le |B|\le n-1$ admits
\[
 f(x)=g(h(x_B),x_{B^c})
\]
with Boolean $h$ and $g$. Constants and the two unary literals are
terminal objects, not members of the prime class. Every essential binary
gate is prime. Negations may be absorbed into adjacent gates; input
permutations and input/output negations preserve primality.

\begin{lemma}[Exact criterion]
For essential $f$, such a block $B$ admits a factorization if and only if
the rows $r_a=(f(a,b):b\in\{-1,1\}^{B^c})$, indexed by block assignments
$a$, have exactly two distinct values.
\end{lemma}
\begin{proof}
A one-bit output $h(a)$ permits at most two rows. Essentiality excludes
one row, since that would make every coordinate in $B$ irrelevant.
Conversely label the two rows by signs, let $h(a)$ be their label, and
let $g(s,b)$ be the $b$-entry of the row labeled $s$. This gives the
identity pointwise. $h$ is nonconstant and hence surjective. Each input
of $h$ is essential, because otherwise the same input of $f$ would be
irrelevant. The new wire of $g$ is essential because the rows differ.
Every other input of $g$ is essential by essentiality of $f$ and
surjectivity of $h$.
\end{proof}

\begin{theorem}[Existence and termination]
Every Boolean function has a formula with each relevant physical input
appearing once, whose internal gates are essential primes and whose
terminal operations are constants or literals. No uniqueness is claimed.
\end{theorem}
\begin{proof}
Induct on the number of essential inputs. Constants and unary functions
are terminal. If $f$ is prime stop. Otherwise use the criterion to split
off $h$ on $k=|B|$ inputs and $g$ on $n-k+1$ inputs. Both arities are
strictly below $n$. Apply induction to each, and replace the fresh
literal in the formula for $g$ by the formula for $h$. All other leaves
of $g$ use $B^c$; the $h$ leaves use $B$, so physical supports remain
disjoint. This also proves termination. After unary operations are
absorbed, a tree with $n$ physical leaves and internal fan-in at least
two has at most $n-1$ internal gates and depth at most $n-1$.
\end{proof}

This is equivalent to irreducibility under arbitrary proper disjoint
substitution: any substitution with at least two nonconstant children
and at least one child on two or more essential coordinates gives a
block as above by expanding the other children into the outer gate.
Unary outer gates, singleton children alone, and constants are trivial
factorizations and are excluded. Shared physical variables are excluded.

Signed parity of arity at least three splits off any pair of inputs,
so only binary parity is prime. AND and OR of arity at least three
likewise split off pairs. In contrast the symmetric threshold
$1[\sum_i y_i\ge t]$ on bits is prime for every $2\le t\le n-1$.
Indeed a block of size $k$ gives thresholds $t-j$ on the other $d=n-k$
bits, $j=0,\ldots,k$. Distinct rows are indexed by the thresholds
clipped to $0,\ldots,d+1$. The interval has endpoints
$L=\max(0,t-k)$ and $U=\min(d+1,t)$, and $U-L\ge2$: when $L=0$,
$U\ge2$; otherwise either $U=t$ gives difference $k\ge2$, or
$U=d+1$ gives difference $n+1-t\ge2$. Thus every candidate block
has at least three rows. In particular the prime class has unbounded
arity; factoring alone does not make the remaining gate bound easier.
\end{document}
